\clearpage
\fancyhead[L]{\small BEER UND MÄDLER.}
\hypertarget{native-657-lower}{}
\begin{center}\small Göttingische gelehrte Anzeigen. Stück 137. Seiten 1361..1365. 30. August 1841.\\[3pt]\rule{.72\linewidth}{.3pt}\end{center}

\emph{Beiträge zur physischen Kenntnis der himmlischen Körper im Sonnensystem von} \textsc{Wilhelm Beer}, \emph{königlich-preußischem Geheimen Kommerzienrat, und Dr.} J. H. \textsc{Mädler}, \emph{kaiserlich-russischem Hofrat.} Weimar 1841. Bei B. F. Voigt. 152 Seiten im Quartformat, dazu 7 Kupfertafeln.

\emph{Fragments sur les corps célestes du système solaire par} \textsc{Guillaume Beer} \emph{et} \textsc{Jean Henri Mädler.} Paris 1840. Chez Bachelier. 216 Seiten im Quartformat, mit 7 Tafeln.

Aus denselben Händen, denen wir die vor vier Jahren erschienene klassische Selenografie verdanken, erhalten wir hier in deutscher und französischer Sprache eine Sammlung kurzer Aufsätze, von denen einige schon früher teilweise bekannt gemacht worden waren. Sie beziehen sich vor allem auf die physische Beschaffenheit des Mondes und der Planeten. Sie zeugen von sorgfältiger Beobachtung, strenger Prüfung und gesundem Urteil und erweitern unsere Kenntnisse in mancher Hinsicht auf wertvolle Weise. Die einzelnen Aufsätze behandeln folgende Gegenstände.
I. \emph{Über die erdabgewandte Mondhalbkugel.} Die von uns abgewandte Hälfte der Mondoberfläche zerfällt in zwei Teile, einen, der bei günstiger Libration für uns sichtbar wird, und einen anderen, der von der Erde aus niemals gesehen werden kann. Nach der hier gegebenen Beschreibung wird dieser zweite Teil von einer Ellipse mit doppelter seitlicher Abplattung begrenzt, und seine Fläche wird mit 0.4243 der gesamten Mondoberfläche angegeben. Der Rezensent findet, dass die Form ein \emph{Viereck} ist und dessen Flächeninhalt 0.4118 beträgt. Über die physische Beschaffenheit dieses Teils können wir zwar niemals gesicherte Kenntnis erlangen. Die Verfasser bemerken aber mit Recht, dass es keinerlei Grund gibt, in dieser Hinsicht einen wesentlichen Unterschied zwischen den beiden Hälften des Mondes zu vermuten, zumal der durch die Libration sichtbar werdende Teil der erdabgewandten Hälfte ganz ähnliche Formationen zeigt wie die erdzugewandte Hälfte. Der einzige Unterschied zwischen den beiden Hälften, den wir mit Gewissheii kennen, liegt in ihren unterschiedlichen Beleuchtungsverhältnissen. Die Verfasser legen diese sehr ausführlich dar, bis hin zu den Vorteilen für die beobachtende Astronomie, die die erdabgewandte Mondhalbkugel gegenüber anderen Standorten hätte, wenn sich der Astronom in Gedanken dorthin versetzt.
II. \emph{Über die Rillen in der Mondoberfläche.} Nach \textsc{Schröter}, der vor 50 \SourceLeaf{658}Jahren als Erster die beiden Rillen bei Hyginus und Ariadäus entdeckt hat, wurde noch eine Reihe weiterer von \textsc{Pastorf}, \textsc{Gruithuysen} und besonders von \textsc{Lohrmann} beobachtet. Ihre Gesamtzahl ist aber unbedeutend im Vergleich mit der großen Menge ähnlicher Formationen, die die Verfasser bei der Arbeit an ihrer Mondkarte gefunden haben. Das hier mitgeteilte Verzeichnis enthält 92 davon.
III. \emph{Über Mondfinsternisse.} Formeln zu ihrer Berechnung, besonders in Bezug auf die Vergrößerung, die man am Radius des Schattenkegels vornehmen muss, um die Beobachtungen mit der Rechnung in Übereinstimmung zu bringen. Die Anwendung auf die drei Finsternisse vom 26. Dezember 1833, 10. Juni 1835 und 13. Oktober 1837 ergibt sehr unterschiedliche Resultate, die erste $\frac{1}{65.4}$, die zweite $\frac{1}{28.3}$, die dritte $\frac{1}{54.0}$. Als Ursache dieser Erscheinung betrachtet man gewöhnlich die Erdatmosphäre, die nur einen geringen Teil der Sonnenstrahlen durchlässt und so den Durchmesser des undurchsichtigen Erdkörpers und damit auch den des Schattenkegels vergrößert. Diese Erklärung ist jedoch auf jeden Fall sehr unbefriedigend, wenn man die Erscheinung in quantitativer Hinsicht genauer betrachtet. Eine starke Lichtabsorption lässt sich nur für die unteren, dichteren Luftschichten annehmen, während das durch die höheren Luftschichten hindurchgehende Licht immer weniger abgeschwächt wird und zugleich durch die Refraktion den Durchmesser des Schattenkegels sogar verringert. Um die Erscheinung, wie sie bei der Mondfinsternis vom 10. Juni 1835 beobachtet wurde, aus dieser Ursache zu erklären, müsste man eine undurchsichtige Atmosphäre von 23 Meilen Höhe annehmen. Ohne dieser Ursache also eine Mitwirkung abzusprechen, müssen wir zugeben, dass die Beobachter einer Mondfinsternis einen Punkt der Mondscheibe schon in den vollen Schatten eintreten zu sehen glauben, während er tatsächlich noch von einem kleinen Teil der Sonnenscheibe Licht erhält. Im betreffenden Fall ergibt die Rechnung, dass ohne Berücksichtigung der Abschwächung durch die Erdatmosphäre von den Punkten an der scheinbaren Grenze des Kernschattens aus noch $\frac{1}{43}$ der Sonnenscheibe unbedeckt blieb. Unter derselben Einschränkung ergibt sich $\frac{1}{178}$ für die Finsternis von 1833 und $\frac{1}{126}$ für die von 1837.
Die Aufsätze IV und V enthalten wertvolle Monografien über interessante Mondlandschaften, die Gegend des Kraters \textsc{Schröter} und die Umgebung des Mondnordpols. Der Aufsatz V bezieht sich auf Saturn und enthält einen sehr verdienstvollen Versuch, die Bahnen der beiden inneren Monde aus \textsc{Herschel}s neu ausgewerteten Beobachtungen von 1789 zu bestimmen, eine Darstellung der Erscheinungen des Saturnrings vom Planeten aus gesehen und eine Bemerkung über den grauen Äquatorstreifen auf Saturn. Die Verfasser deuten die Vermutung an, dass es sich dabei vielleicht um eine wasserähnliche Flüssigkeit handelt, die mit der Anziehungskraft des Rings in Verbindung steht. Doch der Behauptung, der Ring müsse in den Gegenden, in denen er im Zenit steht, \emph{notwendig} eine viele tausend Mal größere konstante Flut als auf der Erde hervorrufen, können wir nicht zustimmen. Ob es eine solche konstante Flut gibt, hängt in der Tat wesentlich von der Gestalt des festen Teils des Planeten selbst ab. Hat sich diese Gestalt, wie wir nach Analogie für wahrscheinlich halten müssen, unter dem Einfluss aller zum Saturnsystem gehörenden Kräfte gebildet, also auch derjenigen, die der Anziehung des Rings zuzuschreiben sind, so kann eigentlich von einer konstanten Flut gar keine Rede sein.
VII. Jupiter. Neue Bestimmung seiner Rotationszeit, relative Helligkeit der Monde, Messungen des Äquatorial- und Polardurchmessers.
VIII. Mars, besonders seine weißen Polarflecken.
IX. Venus. Achtzehn Zeichnungen der Gestalt der Lichtgrenze zeigen eigentlich nur, wie schwer es ist, etwas Genaues über die Rotationszeit festzustellen, bei der die Angaben von \textsc{Cassini} und \textsc{Bianchini} so enorm voneinander abweichen. Die Verfasser finden jedoch in diesen Beobachtungen, so wenig sie auch zur Entscheidung geeignet sind, eher einen Hinweis auf die kürzere Rotationszeit. Über eine eigentümliche strahlende Erscheinung, die Herr \textsc{Mädler} am 7. April 1833 an der Venus bemerkte, lässt sich auf der \SourceLeaf{659}Grundlage des Mitgeteilten kein Urteil fällen.
X. Merkur. Messung seines Durchmessers bei seinem Transit vor der Sonne am 4. Mai 1832.
XI. Heliometrische Messungen an einigen Doppelsternen.
